\originalchapter{8.05 作业4：算符指数与正交逼近}
来源为Fall 2013 Assignment 4, 2013年9月27日发布, 10月4日15时截止. 原题第1题署Based on Griffiths, 第7题署based on Axler. 下文用MIT实际题面的算符恒等式与自定函数逼近任务中文重述, 保留这些参考信用, 不取得外书. 解答均为\textbf{AI独立解答}; 第6题只有Shankar 1.8.8、p.43引用, 不补造题面.
\originalsection{第1题：对易子的导数规则}
\begin{exercise}
证乘积对易式与Jacobi恒等式, 从$[q,p]=\ii\hbar$求$[q^n,p]$与可展开函数$[f(q),p]$, 再在位置表示核验.
\end{exercise}
\begin{solution}
展开$[A,B]C+B[A,C]$中间$BAC$抵消, 给$[A,BC]$. 三个$[[A,B],C]$循环项展开后十二项两两抵消. 归纳得$[q^n,p]=\ii\hbar nq^{n-1}$; 在收敛且共同定义域允许逐项运算的级数上,$[f(q),p]=\ii\hbar f'(q)$. 位置作用$[f(x),-\ii\hbar\partial_x]\psi=\ii\hbar f'(x)\psi$, 同号. 数学微分恒等式独立核算, 不把教材引用误读为缺失原件.
\end{solution}
\originalsection{第2题：指数微分和位移}
\begin{exercise}
证明$\partial_t\e^{t(A+B)}=(A+B)\e^{t(A+B)}=\e^{t(A+B)}(A+B)$. 若$[A,B]=cI$, 求$\e^AB\e^{-A}$. 对$T(a)=\e^{-\ii ap/\hbar}$, 求$T^\dagger xT$和位移后的波函数.
\end{exercise}
\begin{solution}
单一矩阵$K=A+B$的幂总与$K$对易, 逐项微分指数便得第一式, 无需$A,B$彼此对易. 令$F(t)=\e^{tA}B\e^{-tA}$, $F'=\e^{tA}[A,B]\e^{-tA}=cI$, 所以$F(1)=B+cI$. 取$A=\ii ap/\hbar$及$B=x$, $[A,x]=aI$, 得$T^\dagger xT=x+a$. 位置表示$T=\e^{-a\partial_x}$, Taylor平移$T\psi(x)=\psi(x-a)$. 密度向右移$a$, 与均值加$a$一致.
\end{solution}
\originalsection{第3题：Hadamard展开}
\begin{exercise}
从$F(t)=\e^{tA}B\e^{-tA}$的各阶导数证$\e^AB\e^{-A}=\sum_{n\ge0}(\operatorname{ad}_A)^nB/n!$, $\operatorname{ad}_A(X)=[A,X]$.
\end{exercise}
\begin{solution}
$F'=[A,F]$, 继续求导得$F''=[A,[A,F]]$, $F'''=[A,[A,[A,F]]]$. 因$A$不变, 归纳$F^{(n)}=\operatorname{ad}_A^nF$, $F^{(n)}(0)=\operatorname{ad}_A^nB$. Taylor在$t=1$给结论. 有限维矩阵级数绝对收敛; 无界算符只能在保证各嵌套作用与收敛的态上用, 不能只凭形式级数跨越定义域.
\end{solution}
\originalsection{第4题：中心对易子的BCH}
\begin{exercise}
$[A,B]=cI$, 用$G(t)=\e^{t(A+B)}\e^{-tA}$证$G^{-1}G'=B+ct$并求解, 推$\e^{A+B}=\e^B\e^A\e^{c/2}=\e^A\e^B\e^{-c/2}$.
\end{exercise}
\begin{solution}
直接求导$G'=(A+B)G-GA$. 左乘逆得到$G^{-1}G'=\e^{tA}(A+B)\e^{-tA}-A=B+ctI$, 用上一题中心对易特例. 因$B$与$cI$对易,$G(t)=\e^{tB}\e^{ct^2/2}$且$G(0)=I$. $t=1$再右乘$\e^A$给第一式; 交换$A,B$使$c$变$-c$给第二式. 相位因子来自不对易, 不能将其忽略为一般的指数相加.
\end{solution}
\originalsection{第5题：bra、外积与秩}
\begin{exercise}
正交基$\ket1,\ket2,\ket3$中$\ket\psi=a\ket1-b\ket2+a\ket3$, $\ket\phi=b\ket1+a\ket2$. 求bra、相互内积与$A=\ket\phi\bra\psi$矩阵; $Q=\ket\psi\bra\psi+\ket\phi\bra\phi$是否Hermitian, 为何有零本征值?
\end{exercise}
\begin{solution}
bra行分别$(a^*,-b^*,a^*)$及$(b^*,a^*,0)$, 所以$\braket\phi\psi=b^*a-a^*b$, 逆内积为其共轭. ket列为$(a,-b,a)^{\mathsf T}$与$(b,a,0)^{\mathsf T}$, 外积
\[
A=\begin{pmatrix}ba^*&-|b|^2&ba^*\\|a|^2&-ab^*&|a|^2\\0&0&0\end{pmatrix}.
\]
每个$\ket v\bra v$都等于自己的伴随且半正定, 故$Q$Hermitian. 它的像包含于$\operatorname{span}\{\psi,\phi\}$, 秩最多2, 三维空间的核至少1维; 与两矢都正交的非零矢被$Q$消去.
\end{solution}
\originalsection{第7题：投影与五次逼近}
\begin{exercise}
(a)证子空间$U$中的最佳近似为$P_Uv$. (b)$[-\pi,\pi]$实函数内积$\int fg$, 对$U=\operatorname{span}(1,x,\ldots,x^5)$求正交基. (c)求$\sin x,\cos x$在$U$的最佳多项式（精确与六位数）. (d)同图比较原函数、投影与Taylor截断.
\end{exercise}
\begin{solution}
$v-u=(v-P_Uv)+(P_Uv-u)$两项正交, 故$\norm{v-u}^2=\norm{v-P_Uv}^2+\norm{P_Uv-u}^2$, 等号只在$u=P_Uv$. Gram--Schmidt给缩放Legendre多项式$e_{n+1}(x)=\sqrt{(2n+1)/(2\pi)}P_n(x/\pi)$,$0\le n\le5$, 其中$P_n(t)$依次为$1,t,(3t^2-1)/2,(5t^3-3t)/2,(35t^4-30t^2+3)/8,(63t^5-70t^3+15t)/8$. 从幂逐一减正交投影恰产生这些式, $\int_{-1}^1P_n^2=2/(2n+1)$给归一化.

令$p_s=\sum_{j=0}^2a_jx^{2j+1}$, $p_c=\sum_{j=0}^2b_jx^{2j}$. 为避免长的展开式遮住来源, 给精确线性方程作为无歧义系数定义:
\[
\sum_{j=0}^2\frac{2\pi^{2(i+j)+3}}{2(i+j)+3}a_j
=\int_{-\pi}^{\pi}x^{2i+1}\sin x\dd x,\qquad
\sum_{j=0}^2\frac{2\pi^{2(i+j)+1}}{2(i+j)+1}b_j
=\int_{-\pi}^{\pi}x^{2i}\cos x\dd x.
\]
两右端向量分别$(2\pi,2\pi^3-12\pi,2\pi^5-40\pi^3+240\pi)$与$(0,-4\pi,-8\pi^3+48\pi)$. 它们为正定Gram矩阵, 逆矩阵乘这些明确的$\pi$多项式就是精确答案. 逐次解Gram系统并整理为
\begin{align*}
p_s={}&\frac{105(\pi^4-153\pi^2+1485)}{8\pi^6}x
-\frac{315(\pi^4-125\pi^2+1155)}{4\pi^8}x^3\\
&+\frac{693(\pi^4-105\pi^2+945)}{8\pi^{10}}x^5,\\
p_c={}&-\frac{105(2\pi^2-27)}{8\pi^4}
+\frac{315(4\pi^2-45)}{4\pi^6}x^2
-\frac{1575(2\pi^2-21)}{8\pi^8}x^4.
\end{align*}
六位有效数字为$p_s=0.987862x-0.155271x^3+0.00564312x^5$,
$p_c=0.978326-0.452288x^2+0.0261598x^4$. 比较图如下; 每个系数由上面线性系统产生, 非手调拟合. Taylor为$x-x^3/6+x^5/120$与$1-x^2/2+x^4/24$. 投影优化整个区间的平方误差, Taylor只匹配原点导数.
\begin{center}\begin{tikzpicture}
\begin{axis}[width=.82\textwidth,height=4.2cm,xlabel={$x$},ylabel={$\sin x$},grid=major,legend style={font=\small},legend pos=north west]
\addplot+[mark=none] coordinates{(-3.14159265,-1.2246468e-16) (-3.1206487,-0.0209424199) (-3.09970475,-0.0418756537) (-3.0787608,-0.0627905195) (-3.05781685,-0.0836778433) (-3.0368729,-0.104528463) (-3.01592895,-0.125333234) (-2.994985,-0.146083029) (-2.97404105,-0.166768747) (-2.95309709,-0.187381315) (-2.93215314,-0.207911691) (-2.91120919,-0.22835087) (-2.89026524,-0.248689887) (-2.86932129,-0.268919821) (-2.84837734,-0.289031797) (-2.82743339,-0.309016994) (-2.80648944,-0.328866647) (-2.78554549,-0.348572047) (-2.76460154,-0.368124553) (-2.74365758,-0.387515586) (-2.72271363,-0.406736643) (-2.70176968,-0.425779292) (-2.68082573,-0.444635179) (-2.65988178,-0.463296035) (-2.63893783,-0.481753674) (-2.61799388,-0.5) (-2.59704993,-0.518027009) (-2.57610598,-0.535826795) (-2.55516202,-0.553391549) (-2.53421807,-0.570713568) (-2.51327412,-0.587785252) (-2.49233017,-0.604599115) (-2.47138622,-0.62114778) (-2.45044227,-0.63742399) (-2.42949832,-0.653420604) (-2.40855437,-0.669130606) (-2.38761042,-0.684547106) (-2.36666647,-0.699663341) (-2.34572251,-0.71447268) (-2.32477856,-0.728968627) (-2.30383461,-0.743144825) (-2.28289066,-0.756995056) (-2.26194671,-0.770513243) (-2.24100276,-0.783693457) (-2.22005881,-0.796529918) (-2.19911486,-0.809016994) (-2.17817091,-0.821149209) (-2.15722696,-0.832921241) (-2.136283,-0.844327926) (-2.11533905,-0.85536426) (-2.0943951,-0.866025404) (-2.07345115,-0.87630668) (-2.0525072,-0.886203579) (-2.03156325,-0.89571176) (-2.0106193,-0.904827052) (-1.98967535,-0.913545458) (-1.9687314,-0.921863152) (-1.94778745,-0.929776486) (-1.92684349,-0.937281989) (-1.90589954,-0.94437637) (-1.88495559,-0.951056516) (-1.86401164,-0.957319498) (-1.84306769,-0.963162567) (-1.82212374,-0.968583161) (-1.80117979,-0.973578903) (-1.78023584,-0.978147601) (-1.75929189,-0.982287251) (-1.73834793,-0.985996037) (-1.71740398,-0.989272333) (-1.69646003,-0.992114701) (-1.67551608,-0.994521895) (-1.65457213,-0.996492859) (-1.63362818,-0.998026728) (-1.61268423,-0.99912283) (-1.59174028,-0.999780683) (-1.57079633,-1) (-1.54985238,-0.999780683) (-1.52890842,-0.99912283) (-1.50796447,-0.998026728) (-1.48702052,-0.996492859) (-1.46607657,-0.994521895) (-1.44513262,-0.992114701) (-1.42418867,-0.989272333) (-1.40324472,-0.985996037) (-1.38230077,-0.982287251) (-1.36135682,-0.978147601) (-1.34041287,-0.973578903) (-1.31946891,-0.968583161) (-1.29852496,-0.963162567) (-1.27758101,-0.957319498) (-1.25663706,-0.951056516) (-1.23569311,-0.94437637) (-1.21474916,-0.937281989) (-1.19380521,-0.929776486) (-1.17286126,-0.921863152) (-1.15191731,-0.913545458) (-1.13097336,-0.904827052) (-1.1100294,-0.89571176) (-1.08908545,-0.886203579) (-1.0681415,-0.87630668) (-1.04719755,-0.866025404) (-1.0262536,-0.85536426) (-1.00530965,-0.844327926) (-0.984365698,-0.832921241) (-0.963421747,-0.821149209) (-0.942477796,-0.809016994) (-0.921533845,-0.796529918) (-0.900589894,-0.783693457) (-0.879645943,-0.770513243) (-0.858701992,-0.756995056) (-0.837758041,-0.743144825) (-0.81681409,-0.728968627) (-0.795870139,-0.71447268) (-0.774926188,-0.699663341) (-0.753982237,-0.684547106) (-0.733038286,-0.669130606) (-0.712094335,-0.653420604) (-0.691150384,-0.63742399) (-0.670206433,-0.62114778) (-0.649262482,-0.604599115) (-0.628318531,-0.587785252) (-0.60737458,-0.570713568) (-0.586430629,-0.553391549) (-0.565486678,-0.535826795) (-0.544542727,-0.518027009) (-0.523598776,-0.5) (-0.502654825,-0.481753674) (-0.481710874,-0.463296035) (-0.460766923,-0.444635179) (-0.439822972,-0.425779292) (-0.41887902,-0.406736643) (-0.397935069,-0.387515586) (-0.376991118,-0.368124553) (-0.356047167,-0.348572047) (-0.335103216,-0.328866647) (-0.314159265,-0.309016994) (-0.293215314,-0.289031797) (-0.272271363,-0.268919821) (-0.251327412,-0.248689887) (-0.230383461,-0.22835087) (-0.20943951,-0.207911691) (-0.188495559,-0.187381315) (-0.167551608,-0.166768747) (-0.146607657,-0.146083029) (-0.125663706,-0.125333234) (-0.104719755,-0.104528463) (-0.0837758041,-0.0836778433) (-0.0628318531,-0.0627905195) (-0.041887902,-0.0418756537) (-0.020943951,-0.0209424199) (-4.4408921e-16,-4.4408921e-16) (0.020943951,0.0209424199) (0.041887902,0.0418756537) (0.0628318531,0.0627905195) (0.0837758041,0.0836778433) (0.104719755,0.104528463) (0.125663706,0.125333234) (0.146607657,0.146083029) (0.167551608,0.166768747) (0.188495559,0.187381315) (0.20943951,0.207911691) (0.230383461,0.22835087) (0.251327412,0.248689887) (0.272271363,0.268919821) (0.293215314,0.289031797) (0.314159265,0.309016994) (0.335103216,0.328866647) (0.356047167,0.348572047) (0.376991118,0.368124553) (0.397935069,0.387515586) (0.41887902,0.406736643) (0.439822972,0.425779292) (0.460766923,0.444635179) (0.481710874,0.463296035) (0.502654825,0.481753674) (0.523598776,0.5) (0.544542727,0.518027009) (0.565486678,0.535826795) (0.586430629,0.553391549) (0.60737458,0.570713568) (0.628318531,0.587785252) (0.649262482,0.604599115) (0.670206433,0.62114778) (0.691150384,0.63742399) (0.712094335,0.653420604) (0.733038286,0.669130606) (0.753982237,0.684547106) (0.774926188,0.699663341) (0.795870139,0.71447268) (0.81681409,0.728968627) (0.837758041,0.743144825) (0.858701992,0.756995056) (0.879645943,0.770513243) (0.900589894,0.783693457) (0.921533845,0.796529918) (0.942477796,0.809016994) (0.963421747,0.821149209) (0.984365698,0.832921241) (1.00530965,0.844327926) (1.0262536,0.85536426) (1.04719755,0.866025404) (1.0681415,0.87630668) (1.08908545,0.886203579) (1.1100294,0.89571176) (1.13097336,0.904827052) (1.15191731,0.913545458) (1.17286126,0.921863152) (1.19380521,0.929776486) (1.21474916,0.937281989) (1.23569311,0.94437637) (1.25663706,0.951056516) (1.27758101,0.957319498) (1.29852496,0.963162567) (1.31946891,0.968583161) (1.34041287,0.973578903) (1.36135682,0.978147601) (1.38230077,0.982287251) (1.40324472,0.985996037) (1.42418867,0.989272333) (1.44513262,0.992114701) (1.46607657,0.994521895) (1.48702052,0.996492859) (1.50796447,0.998026728) (1.52890842,0.99912283) (1.54985238,0.999780683) (1.57079633,1) (1.59174028,0.999780683) (1.61268423,0.99912283) (1.63362818,0.998026728) (1.65457213,0.996492859) (1.67551608,0.994521895) (1.69646003,0.992114701) (1.71740398,0.989272333) (1.73834793,0.985996037) (1.75929189,0.982287251) (1.78023584,0.978147601) (1.80117979,0.973578903) (1.82212374,0.968583161) (1.84306769,0.963162567) (1.86401164,0.957319498) (1.88495559,0.951056516) (1.90589954,0.94437637) (1.92684349,0.937281989) (1.94778745,0.929776486) (1.9687314,0.921863152) (1.98967535,0.913545458) (2.0106193,0.904827052) (2.03156325,0.89571176) (2.0525072,0.886203579) (2.07345115,0.87630668) (2.0943951,0.866025404) (2.11533905,0.85536426) (2.136283,0.844327926) (2.15722696,0.832921241) (2.17817091,0.821149209) (2.19911486,0.809016994) (2.22005881,0.796529918) (2.24100276,0.783693457) (2.26194671,0.770513243) (2.28289066,0.756995056) (2.30383461,0.743144825) (2.32477856,0.728968627) (2.34572251,0.71447268) (2.36666647,0.699663341) (2.38761042,0.684547106) (2.40855437,0.669130606) (2.42949832,0.653420604) (2.45044227,0.63742399) (2.47138622,0.62114778) (2.49233017,0.604599115) (2.51327412,0.587785252) (2.53421807,0.570713568) (2.55516202,0.553391549) (2.57610598,0.535826795) (2.59704993,0.518027009) (2.61799388,0.5) (2.63893783,0.481753674) (2.65988178,0.463296035) (2.68082573,0.444635179) (2.70176968,0.425779292) (2.72271363,0.406736643) (2.74365758,0.387515586) (2.76460154,0.368124553) (2.78554549,0.348572047) (2.80648944,0.328866647) (2.82743339,0.309016994) (2.84837734,0.289031797) (2.86932129,0.268919821) (2.89026524,0.248689887) (2.91120919,0.22835087) (2.93215314,0.207911691) (2.95309709,0.187381315) (2.97404105,0.166768747) (2.994985,0.146083029) (3.01592895,0.125333234) (3.0368729,0.104528463) (3.05781685,0.0836778433) (3.0787608,0.0627905195) (3.09970475,0.0418756537) (3.1206487,0.0209424199) (3.14159265,1.2246468e-16)};
\addlegendentry{原函数}
\addplot+[mark=none] coordinates{(-3.14159265,-0.0159772936) (-3.1206487,-0.0341337393) (-3.09970475,-0.0525194715) (-3.0787608,-0.0711129591) (-3.05781685,-0.089893075) (-3.0368729,-0.108839093) (-3.01592895,-0.127930686) (-2.994985,-0.147147922) (-2.97404105,-0.166471261) (-2.95309709,-0.185881556) (-2.93215314,-0.205360045) (-2.91120919,-0.224888352) (-2.89026524,-0.244448482) (-2.86932129,-0.26402282) (-2.84837734,-0.283594129) (-2.82743339,-0.303145543) (-2.80648944,-0.322660569) (-2.78554549,-0.342123082) (-2.76460154,-0.361517323) (-2.74365758,-0.380827896) (-2.72271363,-0.400039763) (-2.70176968,-0.419138247) (-2.68082573,-0.438109024) (-2.65988178,-0.456938121) (-2.63893783,-0.475611917) (-2.61799388,-0.494117135) (-2.59704993,-0.512440843) (-2.57610598,-0.53057045) (-2.55516202,-0.548493703) (-2.53421807,-0.566198686) (-2.51327412,-0.583673814) (-2.49233017,-0.600907834) (-2.47138622,-0.617889818) (-2.45044227,-0.634609167) (-2.42949832,-0.651055599) (-2.40855437,-0.667219155) (-2.38761042,-0.68309019) (-2.36666647,-0.698659376) (-2.34572251,-0.713917692) (-2.32477856,-0.728856429) (-2.30383461,-0.743467181) (-2.28289066,-0.757741848) (-2.26194671,-0.771672626) (-2.24100276,-0.785252012) (-2.22005881,-0.798472796) (-2.19911486,-0.811328061) (-2.17817091,-0.823811177) (-2.15722696,-0.835915805) (-2.136283,-0.847635885) (-2.11533905,-0.85896564) (-2.0943951,-0.869899573) (-2.07345115,-0.88043246) (-2.0525072,-0.890559351) (-2.03156325,-0.900275567) (-2.0106193,-0.909576695) (-1.98967535,-0.918458587) (-1.9687314,-0.926917357) (-1.94778745,-0.93494938) (-1.92684349,-0.942551285) (-1.90589954,-0.949719955) (-1.88495559,-0.956452527) (-1.86401164,-0.962746383) (-1.84306769,-0.968599152) (-1.82212374,-0.974008705) (-1.80117979,-0.978973155) (-1.78023584,-0.983490851) (-1.75929189,-0.987560377) (-1.73834793,-0.991180548) (-1.71740398,-0.99435041) (-1.69646003,-0.997069235) (-1.67551608,-0.999336518) (-1.65457213,-1.00115197) (-1.63362818,-1.00251554) (-1.61268423,-1.00342737) (-1.59174028,-1.00388782) (-1.57079633,-1.00389746) (-1.54985238,-1.00345708) (-1.52890842,-1.00256767) (-1.50796447,-1.0012304) (-1.48702052,-0.999446672) (-1.46607657,-0.997218062) (-1.44513262,-0.994546349) (-1.42418867,-0.991433499) (-1.40324472,-0.987881669) (-1.38230077,-0.983893202) (-1.36135682,-0.979470621) (-1.34041287,-0.974616631) (-1.31946891,-0.969334115) (-1.29852496,-0.963626128) (-1.27758101,-0.9574959) (-1.25663706,-0.950946828) (-1.23569311,-0.943982477) (-1.21474916,-0.936606574) (-1.19380521,-0.928823009) (-1.17286126,-0.920635829) (-1.15191731,-0.912049237) (-1.13097336,-0.903067588) (-1.1100294,-0.893695388) (-1.08908545,-0.883937291) (-1.0681415,-0.873798094) (-1.04719755,-0.863282737) (-1.0262536,-0.852396298) (-1.00530965,-0.841143993) (-0.984365698,-0.829531171) (-0.963421747,-0.817563311) (-0.942477796,-0.805246023) (-0.921533845,-0.79258504) (-0.900589894,-0.779586218) (-0.879645943,-0.766255535) (-0.858701992,-0.752599083) (-0.837758041,-0.738623073) (-0.81681409,-0.724333823) (-0.795870139,-0.709737765) (-0.774926188,-0.694841433) (-0.753982237,-0.679651468) (-0.733038286,-0.664174611) (-0.712094335,-0.6484177) (-0.691150384,-0.632387671) (-0.670206433,-0.616091549) (-0.649262482,-0.599536454) (-0.628318531,-0.582729589) (-0.60737458,-0.565678243) (-0.586430629,-0.548389788) (-0.565486678,-0.530871675) (-0.544542727,-0.513131429) (-0.523598776,-0.495176651) (-0.502654825,-0.477015013) (-0.481710874,-0.458654254) (-0.460766923,-0.440102179) (-0.439822972,-0.421366655) (-0.41887902,-0.402455612) (-0.397935069,-0.383377033) (-0.376991118,-0.364138959) (-0.356047167,-0.344749481) (-0.335103216,-0.32521674) (-0.314159265,-0.305548924) (-0.293215314,-0.285754261) (-0.272271363,-0.265841026) (-0.251327412,-0.245817526) (-0.230383461,-0.225692108) (-0.20943951,-0.205473151) (-0.188495559,-0.185169061) (-0.167551608,-0.164788274) (-0.146607657,-0.144339251) (-0.125663706,-0.123830473) (-0.104719755,-0.103270441) (-0.0837758041,-0.082667673) (-0.0628318531,-0.062030699) (-0.041887902,-0.0413680612) (-0.020943951,-0.0206883097) (-4.4408921e-16,-4.38698915e-16) (0.020943951,0.0206883097) (0.041887902,0.0413680612) (0.0628318531,0.062030699) (0.0837758041,0.082667673) (0.104719755,0.103270441) (0.125663706,0.123830473) (0.146607657,0.144339251) (0.167551608,0.164788274) (0.188495559,0.185169061) (0.20943951,0.205473151) (0.230383461,0.225692108) (0.251327412,0.245817526) (0.272271363,0.265841026) (0.293215314,0.285754261) (0.314159265,0.305548924) (0.335103216,0.32521674) (0.356047167,0.344749481) (0.376991118,0.364138959) (0.397935069,0.383377033) (0.41887902,0.402455612) (0.439822972,0.421366655) (0.460766923,0.440102179) (0.481710874,0.458654254) (0.502654825,0.477015013) (0.523598776,0.495176651) (0.544542727,0.513131429) (0.565486678,0.530871675) (0.586430629,0.548389788) (0.60737458,0.565678243) (0.628318531,0.582729589) (0.649262482,0.599536454) (0.670206433,0.616091549) (0.691150384,0.632387671) (0.712094335,0.6484177) (0.733038286,0.664174611) (0.753982237,0.679651468) (0.774926188,0.694841433) (0.795870139,0.709737765) (0.81681409,0.724333823) (0.837758041,0.738623073) (0.858701992,0.752599083) (0.879645943,0.766255535) (0.900589894,0.779586218) (0.921533845,0.79258504) (0.942477796,0.805246023) (0.963421747,0.817563311) (0.984365698,0.829531171) (1.00530965,0.841143993) (1.0262536,0.852396298) (1.04719755,0.863282737) (1.0681415,0.873798094) (1.08908545,0.883937291) (1.1100294,0.893695388) (1.13097336,0.903067588) (1.15191731,0.912049237) (1.17286126,0.920635829) (1.19380521,0.928823009) (1.21474916,0.936606574) (1.23569311,0.943982477) (1.25663706,0.950946828) (1.27758101,0.9574959) (1.29852496,0.963626128) (1.31946891,0.969334115) (1.34041287,0.974616631) (1.36135682,0.979470621) (1.38230077,0.983893202) (1.40324472,0.987881669) (1.42418867,0.991433499) (1.44513262,0.994546349) (1.46607657,0.997218062) (1.48702052,0.999446672) (1.50796447,1.0012304) (1.52890842,1.00256767) (1.54985238,1.00345708) (1.57079633,1.00389746) (1.59174028,1.00388782) (1.61268423,1.00342737) (1.63362818,1.00251554) (1.65457213,1.00115197) (1.67551608,0.999336518) (1.69646003,0.997069235) (1.71740398,0.99435041) (1.73834793,0.991180548) (1.75929189,0.987560377) (1.78023584,0.983490851) (1.80117979,0.978973155) (1.82212374,0.974008705) (1.84306769,0.968599152) (1.86401164,0.962746383) (1.88495559,0.956452527) (1.90589954,0.949719955) (1.92684349,0.942551285) (1.94778745,0.93494938) (1.9687314,0.926917357) (1.98967535,0.918458587) (2.0106193,0.909576695) (2.03156325,0.900275567) (2.0525072,0.890559351) (2.07345115,0.88043246) (2.0943951,0.869899573) (2.11533905,0.85896564) (2.136283,0.847635885) (2.15722696,0.835915805) (2.17817091,0.823811177) (2.19911486,0.811328061) (2.22005881,0.798472796) (2.24100276,0.785252012) (2.26194671,0.771672626) (2.28289066,0.757741848) (2.30383461,0.743467181) (2.32477856,0.728856429) (2.34572251,0.713917692) (2.36666647,0.698659376) (2.38761042,0.68309019) (2.40855437,0.667219155) (2.42949832,0.651055599) (2.45044227,0.634609167) (2.47138622,0.617889818) (2.49233017,0.600907834) (2.51327412,0.583673814) (2.53421807,0.566198686) (2.55516202,0.548493703) (2.57610598,0.53057045) (2.59704993,0.512440843) (2.61799388,0.494117135) (2.63893783,0.475611917) (2.65988178,0.456938121) (2.68082573,0.438109024) (2.70176968,0.419138247) (2.72271363,0.400039763) (2.74365758,0.380827896) (2.76460154,0.361517323) (2.78554549,0.342123082) (2.80648944,0.322660569) (2.82743339,0.303145543) (2.84837734,0.283594129) (2.86932129,0.26402282) (2.89026524,0.244448482) (2.91120919,0.224888352) (2.93215314,0.205360045) (2.95309709,0.185881556) (2.97404105,0.166471261) (2.994985,0.147147922) (3.01592895,0.127930686) (3.0368729,0.108839093) (3.05781685,0.089893075) (3.0787608,0.0711129591) (3.09970475,0.0525194715) (3.1206487,0.0341337393) (3.14159265,0.0159772936)};
\addlegendentry{正交投影}
\addplot+[mark=none] coordinates{(-3.14159265,-0.524043913) (-3.1206487,-0.521887128) (-3.09970475,-0.520583301) (-3.0787608,-0.520097185) (-3.05781685,-0.520394128) (-3.0368729,-0.521440071) (-3.01592895,-0.523201546) (-2.994985,-0.525645665) (-2.97404105,-0.528740122) (-2.95309709,-0.532453189) (-2.93215314,-0.536753709) (-2.91120919,-0.541611091) (-2.89026524,-0.546995312) (-2.86932129,-0.552876907) (-2.84837734,-0.559226967) (-2.82743339,-0.566017135) (-2.80648944,-0.573219605) (-2.78554549,-0.580807111) (-2.76460154,-0.588752928) (-2.74365758,-0.59703087) (-2.72271363,-0.605615278) (-2.70176968,-0.614481026) (-2.68082573,-0.623603507) (-2.65988178,-0.632958638) (-2.63893783,-0.64252285) (-2.61799388,-0.652273085) (-2.59704993,-0.662186794) (-2.57610598,-0.672241932) (-2.55516202,-0.682416953) (-2.53421807,-0.692690805) (-2.51327412,-0.703042932) (-2.49233017,-0.713453262) (-2.47138622,-0.723902207) (-2.45044227,-0.73437066) (-2.42949832,-0.744839988) (-2.40855437,-0.755292031) (-2.38761042,-0.765709095) (-2.36666647,-0.776073951) (-2.34572251,-0.786369827) (-2.32477856,-0.79658041) (-2.30383461,-0.806689834) (-2.28289066,-0.816682685) (-2.26194671,-0.826543987) (-2.24100276,-0.836259209) (-2.22005881,-0.84581425) (-2.19911486,-0.855195444) (-2.17817091,-0.86438955) (-2.15722696,-0.873383751) (-2.136283,-0.88216565) (-2.11533905,-0.890723262) (-2.0943951,-0.899045016) (-2.07345115,-0.907119748) (-2.0525072,-0.914936695) (-2.03156325,-0.922485494) (-2.0106193,-0.929756178) (-1.98967535,-0.936739169) (-1.9687314,-0.943425277) (-1.94778745,-0.949805694) (-1.92684349,-0.955871992) (-1.90589954,-0.961616118) (-1.88495559,-0.967030387) (-1.86401164,-0.972107485) (-1.84306769,-0.976840456) (-1.82212374,-0.981222706) (-1.80117979,-0.985247995) (-1.78023584,-0.988910432) (-1.75929189,-0.992204476) (-1.73834793,-0.995124924) (-1.71740398,-0.997666915) (-1.69646003,-0.999825921) (-1.67551608,-1.00159775) (-1.65457213,-1.00297852) (-1.63362818,-1.00396469) (-1.61268423,-1.00455302) (-1.59174028,-1.00474062) (-1.57079633,-1.00452486) (-1.54985238,-1.00390344) (-1.52890842,-1.00287438) (-1.50796447,-1.00143597) (-1.48702052,-0.999586814) (-1.46607657,-0.997325787) (-1.44513262,-0.994652067) (-1.42418867,-0.991565106) (-1.40324472,-0.988064637) (-1.38230077,-0.984150665) (-1.36135682,-0.979823466) (-1.34041287,-0.975083582) (-1.31946891,-0.969931818) (-1.29852496,-0.964369235) (-1.27758101,-0.958397148) (-1.25663706,-0.952017123) (-1.23569311,-0.945230972) (-1.21474916,-0.938040747) (-1.19380521,-0.93044874) (-1.17286126,-0.922457475) (-1.15191731,-0.914069707) (-1.13097336,-0.905288415) (-1.1100294,-0.896116801) (-1.08908545,-0.886558284) (-1.0681415,-0.876616498) (-1.04719755,-0.866295284) (-1.0262536,-0.855598691) (-1.00530965,-0.844530967) (-0.984365698,-0.833096561) (-0.963421747,-0.821300112) (-0.942477796,-0.80914645) (-0.921533845,-0.796640589) (-0.900589894,-0.783787728) (-0.879645943,-0.770593238) (-0.858701992,-0.757062668) (-0.837758041,-0.743201733) (-0.81681409,-0.729016316) (-0.795870139,-0.714512458) (-0.774926188,-0.69969636) (-0.753982237,-0.684574375) (-0.733038286,-0.669153005) (-0.712094335,-0.653438897) (-0.691150384,-0.637438839) (-0.670206433,-0.621159757) (-0.649262482,-0.604608708) (-0.628318531,-0.587792881) (-0.60737458,-0.570719587) (-0.586430629,-0.553396259) (-0.565486678,-0.535830448) (-0.544542727,-0.518029815) (-0.523598776,-0.500002133) (-0.502654825,-0.481755277) (-0.481710874,-0.463297225) (-0.460766923,-0.444636051) (-0.439822972,-0.425779922) (-0.41887902,-0.406737091) (-0.397935069,-0.387515899) (-0.376991118,-0.368124767) (-0.356047167,-0.348572191) (-0.335103216,-0.328866741) (-0.314159265,-0.309017054) (-0.293215314,-0.289031834) (-0.272271363,-0.268919843) (-0.251327412,-0.2486899) (-0.230383461,-0.228350877) (-0.20943951,-0.207911694) (-0.188495559,-0.187381316) (-0.167551608,-0.166768747) (-0.146607657,-0.146083029) (-0.125663706,-0.125333234) (-0.104719755,-0.104528463) (-0.0837758041,-0.0836778433) (-0.0628318531,-0.0627905195) (-0.041887902,-0.0418756537) (-0.020943951,-0.0209424199) (-4.4408921e-16,-4.4408921e-16) (0.020943951,0.0209424199) (0.041887902,0.0418756537) (0.0628318531,0.0627905195) (0.0837758041,0.0836778433) (0.104719755,0.104528463) (0.125663706,0.125333234) (0.146607657,0.146083029) (0.167551608,0.166768747) (0.188495559,0.187381316) (0.20943951,0.207911694) (0.230383461,0.228350877) (0.251327412,0.2486899) (0.272271363,0.268919843) (0.293215314,0.289031834) (0.314159265,0.309017054) (0.335103216,0.328866741) (0.356047167,0.348572191) (0.376991118,0.368124767) (0.397935069,0.387515899) (0.41887902,0.406737091) (0.439822972,0.425779922) (0.460766923,0.444636051) (0.481710874,0.463297225) (0.502654825,0.481755277) (0.523598776,0.500002133) (0.544542727,0.518029815) (0.565486678,0.535830448) (0.586430629,0.553396259) (0.60737458,0.570719587) (0.628318531,0.587792881) (0.649262482,0.604608708) (0.670206433,0.621159757) (0.691150384,0.637438839) (0.712094335,0.653438897) (0.733038286,0.669153005) (0.753982237,0.684574375) (0.774926188,0.69969636) (0.795870139,0.714512458) (0.81681409,0.729016316) (0.837758041,0.743201733) (0.858701992,0.757062668) (0.879645943,0.770593238) (0.900589894,0.783787728) (0.921533845,0.796640589) (0.942477796,0.80914645) (0.963421747,0.821300112) (0.984365698,0.833096561) (1.00530965,0.844530967) (1.0262536,0.855598691) (1.04719755,0.866295284) (1.0681415,0.876616498) (1.08908545,0.886558284) (1.1100294,0.896116801) (1.13097336,0.905288415) (1.15191731,0.914069707) (1.17286126,0.922457475) (1.19380521,0.93044874) (1.21474916,0.938040747) (1.23569311,0.945230972) (1.25663706,0.952017123) (1.27758101,0.958397148) (1.29852496,0.964369235) (1.31946891,0.969931818) (1.34041287,0.975083582) (1.36135682,0.979823466) (1.38230077,0.984150665) (1.40324472,0.988064637) (1.42418867,0.991565106) (1.44513262,0.994652067) (1.46607657,0.997325787) (1.48702052,0.999586814) (1.50796447,1.00143597) (1.52890842,1.00287438) (1.54985238,1.00390344) (1.57079633,1.00452486) (1.59174028,1.00474062) (1.61268423,1.00455302) (1.63362818,1.00396469) (1.65457213,1.00297852) (1.67551608,1.00159775) (1.69646003,0.999825921) (1.71740398,0.997666915) (1.73834793,0.995124924) (1.75929189,0.992204476) (1.78023584,0.988910432) (1.80117979,0.985247995) (1.82212374,0.981222706) (1.84306769,0.976840456) (1.86401164,0.972107485) (1.88495559,0.967030387) (1.90589954,0.961616118) (1.92684349,0.955871992) (1.94778745,0.949805694) (1.9687314,0.943425277) (1.98967535,0.936739169) (2.0106193,0.929756178) (2.03156325,0.922485494) (2.0525072,0.914936695) (2.07345115,0.907119748) (2.0943951,0.899045016) (2.11533905,0.890723262) (2.136283,0.88216565) (2.15722696,0.873383751) (2.17817091,0.86438955) (2.19911486,0.855195444) (2.22005881,0.84581425) (2.24100276,0.836259209) (2.26194671,0.826543987) (2.28289066,0.816682685) (2.30383461,0.806689834) (2.32477856,0.79658041) (2.34572251,0.786369827) (2.36666647,0.776073951) (2.38761042,0.765709095) (2.40855437,0.755292031) (2.42949832,0.744839988) (2.45044227,0.73437066) (2.47138622,0.723902207) (2.49233017,0.713453262) (2.51327412,0.703042932) (2.53421807,0.692690805) (2.55516202,0.682416953) (2.57610598,0.672241932) (2.59704993,0.662186794) (2.61799388,0.652273085) (2.63893783,0.64252285) (2.65988178,0.632958638) (2.68082573,0.623603507) (2.70176968,0.614481026) (2.72271363,0.605615278) (2.74365758,0.59703087) (2.76460154,0.588752928) (2.78554549,0.580807111) (2.80648944,0.573219605) (2.82743339,0.566017135) (2.84837734,0.559226967) (2.86932129,0.552876907) (2.89026524,0.546995312) (2.91120919,0.541611091) (2.93215314,0.536753709) (2.95309709,0.532453189) (2.97404105,0.528740122) (2.994985,0.525645665) (3.01592895,0.523201546) (3.0368729,0.521440071) (3.05781685,0.520394128) (3.0787608,0.520097185) (3.09970475,0.520583301) (3.1206487,0.521887128) (3.14159265,0.524043913)};
\addlegendentry{Taylor}
\end{axis}\end{tikzpicture}\end{center}
\begin{center}\begin{tikzpicture}
\begin{axis}[width=.82\textwidth,height=4.2cm,xlabel={$x$},ylabel={$\cos x$},grid=major,legend style={font=\small},legend pos=north west]
\addplot+[mark=none] coordinates{(-3.14159265,-1) (-3.1206487,-0.999780683) (-3.09970475,-0.99912283) (-3.0787608,-0.998026728) (-3.05781685,-0.996492859) (-3.0368729,-0.994521895) (-3.01592895,-0.992114701) (-2.994985,-0.989272333) (-2.97404105,-0.985996037) (-2.95309709,-0.982287251) (-2.93215314,-0.978147601) (-2.91120919,-0.973578903) (-2.89026524,-0.968583161) (-2.86932129,-0.963162567) (-2.84837734,-0.957319498) (-2.82743339,-0.951056516) (-2.80648944,-0.94437637) (-2.78554549,-0.937281989) (-2.76460154,-0.929776486) (-2.74365758,-0.921863152) (-2.72271363,-0.913545458) (-2.70176968,-0.904827052) (-2.68082573,-0.89571176) (-2.65988178,-0.886203579) (-2.63893783,-0.87630668) (-2.61799388,-0.866025404) (-2.59704993,-0.85536426) (-2.57610598,-0.844327926) (-2.55516202,-0.832921241) (-2.53421807,-0.821149209) (-2.51327412,-0.809016994) (-2.49233017,-0.796529918) (-2.47138622,-0.783693457) (-2.45044227,-0.770513243) (-2.42949832,-0.756995056) (-2.40855437,-0.743144825) (-2.38761042,-0.728968627) (-2.36666647,-0.71447268) (-2.34572251,-0.699663341) (-2.32477856,-0.684547106) (-2.30383461,-0.669130606) (-2.28289066,-0.653420604) (-2.26194671,-0.63742399) (-2.24100276,-0.62114778) (-2.22005881,-0.604599115) (-2.19911486,-0.587785252) (-2.17817091,-0.570713568) (-2.15722696,-0.553391549) (-2.136283,-0.535826795) (-2.11533905,-0.518027009) (-2.0943951,-0.5) (-2.07345115,-0.481753674) (-2.0525072,-0.463296035) (-2.03156325,-0.444635179) (-2.0106193,-0.425779292) (-1.98967535,-0.406736643) (-1.9687314,-0.387515586) (-1.94778745,-0.368124553) (-1.92684349,-0.348572047) (-1.90589954,-0.328866647) (-1.88495559,-0.309016994) (-1.86401164,-0.289031797) (-1.84306769,-0.268919821) (-1.82212374,-0.248689887) (-1.80117979,-0.22835087) (-1.78023584,-0.207911691) (-1.75929189,-0.187381315) (-1.73834793,-0.166768747) (-1.71740398,-0.146083029) (-1.69646003,-0.125333234) (-1.67551608,-0.104528463) (-1.65457213,-0.0836778433) (-1.63362818,-0.0627905195) (-1.61268423,-0.0418756537) (-1.59174028,-0.0209424199) (-1.57079633,-1.60812265e-16) (-1.54985238,0.0209424199) (-1.52890842,0.0418756537) (-1.50796447,0.0627905195) (-1.48702052,0.0836778433) (-1.46607657,0.104528463) (-1.44513262,0.125333234) (-1.42418867,0.146083029) (-1.40324472,0.166768747) (-1.38230077,0.187381315) (-1.36135682,0.207911691) (-1.34041287,0.22835087) (-1.31946891,0.248689887) (-1.29852496,0.268919821) (-1.27758101,0.289031797) (-1.25663706,0.309016994) (-1.23569311,0.328866647) (-1.21474916,0.348572047) (-1.19380521,0.368124553) (-1.17286126,0.387515586) (-1.15191731,0.406736643) (-1.13097336,0.425779292) (-1.1100294,0.444635179) (-1.08908545,0.463296035) (-1.0681415,0.481753674) (-1.04719755,0.5) (-1.0262536,0.518027009) (-1.00530965,0.535826795) (-0.984365698,0.553391549) (-0.963421747,0.570713568) (-0.942477796,0.587785252) (-0.921533845,0.604599115) (-0.900589894,0.62114778) (-0.879645943,0.63742399) (-0.858701992,0.653420604) (-0.837758041,0.669130606) (-0.81681409,0.684547106) (-0.795870139,0.699663341) (-0.774926188,0.71447268) (-0.753982237,0.728968627) (-0.733038286,0.743144825) (-0.712094335,0.756995056) (-0.691150384,0.770513243) (-0.670206433,0.783693457) (-0.649262482,0.796529918) (-0.628318531,0.809016994) (-0.60737458,0.821149209) (-0.586430629,0.832921241) (-0.565486678,0.844327926) (-0.544542727,0.85536426) (-0.523598776,0.866025404) (-0.502654825,0.87630668) (-0.481710874,0.886203579) (-0.460766923,0.89571176) (-0.439822972,0.904827052) (-0.41887902,0.913545458) (-0.397935069,0.921863152) (-0.376991118,0.929776486) (-0.356047167,0.937281989) (-0.335103216,0.94437637) (-0.314159265,0.951056516) (-0.293215314,0.957319498) (-0.272271363,0.963162567) (-0.251327412,0.968583161) (-0.230383461,0.973578903) (-0.20943951,0.978147601) (-0.188495559,0.982287251) (-0.167551608,0.985996037) (-0.146607657,0.989272333) (-0.125663706,0.992114701) (-0.104719755,0.994521895) (-0.0837758041,0.996492859) (-0.0628318531,0.998026728) (-0.041887902,0.99912283) (-0.020943951,0.999780683) (-4.4408921e-16,1) (0.020943951,0.999780683) (0.041887902,0.99912283) (0.0628318531,0.998026728) (0.0837758041,0.996492859) (0.104719755,0.994521895) (0.125663706,0.992114701) (0.146607657,0.989272333) (0.167551608,0.985996037) (0.188495559,0.982287251) (0.20943951,0.978147601) (0.230383461,0.973578903) (0.251327412,0.968583161) (0.272271363,0.963162567) (0.293215314,0.957319498) (0.314159265,0.951056516) (0.335103216,0.94437637) (0.356047167,0.937281989) (0.376991118,0.929776486) (0.397935069,0.921863152) (0.41887902,0.913545458) (0.439822972,0.904827052) (0.460766923,0.89571176) (0.481710874,0.886203579) (0.502654825,0.87630668) (0.523598776,0.866025404) (0.544542727,0.85536426) (0.565486678,0.844327926) (0.586430629,0.832921241) (0.60737458,0.821149209) (0.628318531,0.809016994) (0.649262482,0.796529918) (0.670206433,0.783693457) (0.691150384,0.770513243) (0.712094335,0.756995056) (0.733038286,0.743144825) (0.753982237,0.728968627) (0.774926188,0.71447268) (0.795870139,0.699663341) (0.81681409,0.684547106) (0.837758041,0.669130606) (0.858701992,0.653420604) (0.879645943,0.63742399) (0.900589894,0.62114778) (0.921533845,0.604599115) (0.942477796,0.587785252) (0.963421747,0.570713568) (0.984365698,0.553391549) (1.00530965,0.535826795) (1.0262536,0.518027009) (1.04719755,0.5) (1.0681415,0.481753674) (1.08908545,0.463296035) (1.1100294,0.444635179) (1.13097336,0.425779292) (1.15191731,0.406736643) (1.17286126,0.387515586) (1.19380521,0.368124553) (1.21474916,0.348572047) (1.23569311,0.328866647) (1.25663706,0.309016994) (1.27758101,0.289031797) (1.29852496,0.268919821) (1.31946891,0.248689887) (1.34041287,0.22835087) (1.36135682,0.207911691) (1.38230077,0.187381315) (1.40324472,0.166768747) (1.42418867,0.146083029) (1.44513262,0.125333234) (1.46607657,0.104528463) (1.48702052,0.0836778433) (1.50796447,0.0627905195) (1.52890842,0.0418756537) (1.54985238,0.0209424199) (1.57079633,6.123234e-17) (1.59174028,-0.0209424199) (1.61268423,-0.0418756537) (1.63362818,-0.0627905195) (1.65457213,-0.0836778433) (1.67551608,-0.104528463) (1.69646003,-0.125333234) (1.71740398,-0.146083029) (1.73834793,-0.166768747) (1.75929189,-0.187381315) (1.78023584,-0.207911691) (1.80117979,-0.22835087) (1.82212374,-0.248689887) (1.84306769,-0.268919821) (1.86401164,-0.289031797) (1.88495559,-0.309016994) (1.90589954,-0.328866647) (1.92684349,-0.348572047) (1.94778745,-0.368124553) (1.9687314,-0.387515586) (1.98967535,-0.406736643) (2.0106193,-0.425779292) (2.03156325,-0.444635179) (2.0525072,-0.463296035) (2.07345115,-0.481753674) (2.0943951,-0.5) (2.11533905,-0.518027009) (2.136283,-0.535826795) (2.15722696,-0.553391549) (2.17817091,-0.570713568) (2.19911486,-0.587785252) (2.22005881,-0.604599115) (2.24100276,-0.62114778) (2.26194671,-0.63742399) (2.28289066,-0.653420604) (2.30383461,-0.669130606) (2.32477856,-0.684547106) (2.34572251,-0.699663341) (2.36666647,-0.71447268) (2.38761042,-0.728968627) (2.40855437,-0.743144825) (2.42949832,-0.756995056) (2.45044227,-0.770513243) (2.47138622,-0.783693457) (2.49233017,-0.796529918) (2.51327412,-0.809016994) (2.53421807,-0.821149209) (2.55516202,-0.832921241) (2.57610598,-0.844327926) (2.59704993,-0.85536426) (2.61799388,-0.866025404) (2.63893783,-0.87630668) (2.65988178,-0.886203579) (2.68082573,-0.89571176) (2.70176968,-0.904827052) (2.72271363,-0.913545458) (2.74365758,-0.921863152) (2.76460154,-0.929776486) (2.78554549,-0.937281989) (2.80648944,-0.94437637) (2.82743339,-0.951056516) (2.84837734,-0.957319498) (2.86932129,-0.963162567) (2.89026524,-0.968583161) (2.91120919,-0.973578903) (2.93215314,-0.978147601) (2.95309709,-0.982287251) (2.97404105,-0.985996037) (2.994985,-0.989272333) (3.01592895,-0.992114701) (3.0368729,-0.994521895) (3.05781685,-0.996492859) (3.0787608,-0.998026728) (3.09970475,-0.99912283) (3.1206487,-0.999780683) (3.14159265,-1)};
\addlegendentry{原函数}
\addplot+[mark=none] coordinates{(-3.14159265,-0.937371052) (-3.1206487,-0.945326366) (-3.09970475,-0.95233748) (-3.0787608,-0.958422332) (-3.05781685,-0.963598741) (-3.0368729,-0.967884405) (-3.01592895,-0.971296901) (-2.994985,-0.973853684) (-2.97404105,-0.97557209) (-2.95309709,-0.976469334) (-2.93215314,-0.97656251) (-2.91120919,-0.97586859) (-2.89026524,-0.974404426) (-2.86932129,-0.97218675) (-2.84837734,-0.969232173) (-2.82743339,-0.965557183) (-2.80648944,-0.961178151) (-2.78554549,-0.956111323) (-2.76460154,-0.950372828) (-2.74365758,-0.943978672) (-2.72271363,-0.93694474) (-2.70176968,-0.929286797) (-2.68082573,-0.921020488) (-2.65988178,-0.912161336) (-2.63893783,-0.902724744) (-2.61799388,-0.892725992) (-2.59704993,-0.882180242) (-2.57610598,-0.871102533) (-2.55516202,-0.859507786) (-2.53421807,-0.847410799) (-2.51327412,-0.834826249) (-2.49233017,-0.821768693) (-2.47138622,-0.808252567) (-2.45044227,-0.794292187) (-2.42949832,-0.779901746) (-2.40855437,-0.765095319) (-2.38761042,-0.749886858) (-2.36666647,-0.734290196) (-2.34572251,-0.718319043) (-2.32477856,-0.701986991) (-2.30383461,-0.685307507) (-2.28289066,-0.668293943) (-2.26194671,-0.650959525) (-2.24100276,-0.63331736) (-2.22005881,-0.615380436) (-2.19911486,-0.597161617) (-2.17817091,-0.578673649) (-2.15722696,-0.559929155) (-2.136283,-0.540940639) (-2.11533905,-0.521720482) (-2.0943951,-0.502280948) (-2.07345115,-0.482634176) (-2.0525072,-0.462792186) (-2.03156325,-0.442766878) (-2.0106193,-0.422570029) (-1.98967535,-0.402213299) (-1.9687314,-0.381708222) (-1.94778745,-0.361066216) (-1.92684349,-0.340298575) (-1.90589954,-0.319416473) (-1.88495559,-0.298430965) (-1.86401164,-0.277352983) (-1.84306769,-0.256193338) (-1.82212374,-0.234962723) (-1.80117979,-0.213671707) (-1.78023584,-0.19233074) (-1.75929189,-0.170950151) (-1.73834793,-0.149540147) (-1.71740398,-0.128110816) (-1.69646003,-0.106672124) (-1.67551608,-0.0852339161) (-1.65457213,-0.0638059175) (-1.63362818,-0.042397732) (-1.61268423,-0.0210188426) (-1.59174028,0.000321388316) (-1.57079633,0.0216137193) (-1.54985238,0.0428490296) (-1.52890842,0.0640183195) (-1.50796447,0.0851127097) (-1.48702052,0.106123442) (-1.46607657,0.127041879) (-1.44513262,0.147859504) (-1.42418867,0.168567922) (-1.40324472,0.189158856) (-1.38230077,0.209624153) (-1.36135682,0.22995578) (-1.34041287,0.250145823) (-1.31946891,0.270186491) (-1.29852496,0.290070113) (-1.27758101,0.309789139) (-1.25663706,0.329336139) (-1.23569311,0.348703804) (-1.21474916,0.367884947) (-1.19380521,0.386872501) (-1.17286126,0.40565952) (-1.15191731,0.424239177) (-1.13097336,0.442604769) (-1.1100294,0.460749712) (-1.08908545,0.478667542) (-1.0681415,0.496351918) (-1.04719755,0.513796618) (-1.0262536,0.530995542) (-1.00530965,0.547942709) (-0.984365698,0.564632261) (-0.963421747,0.581058459) (-0.942477796,0.597215687) (-0.921533845,0.613098447) (-0.900589894,0.628701364) (-0.879645943,0.644019183) (-0.858701992,0.65904677) (-0.837758041,0.673779111) (-0.81681409,0.688211314) (-0.795870139,0.702338606) (-0.774926188,0.716156338) (-0.753982237,0.729659979) (-0.733038286,0.74284512) (-0.712094335,0.755707472) (-0.691150384,0.768242867) (-0.670206433,0.780447258) (-0.649262482,0.79231672) (-0.628318531,0.803847447) (-0.60737458,0.815035754) (-0.586430629,0.825878079) (-0.565486678,0.836370977) (-0.544542727,0.846511126) (-0.523598776,0.856295327) (-0.502654825,0.865720497) (-0.481710874,0.874783677) (-0.460766923,0.883482029) (-0.439822972,0.891812834) (-0.41887902,0.899773496) (-0.397935069,0.907361536) (-0.376991118,0.914574601) (-0.356047167,0.921410455) (-0.335103216,0.927866983) (-0.314159265,0.933942194) (-0.293215314,0.939634213) (-0.272271363,0.94494129) (-0.251327412,0.949861794) (-0.230383461,0.954394214) (-0.20943951,0.958537162) (-0.188495559,0.962289369) (-0.167551608,0.965649687) (-0.146607657,0.96861709) (-0.125663706,0.971190671) (-0.104719755,0.973369646) (-0.0837758041,0.975153349) (-0.0628318531,0.976541238) (-0.041887902,0.977532889) (-0.020943951,0.978128) (-4.4408921e-16,0.978326391) (0.020943951,0.978128) (0.041887902,0.977532889) (0.0628318531,0.976541238) (0.0837758041,0.975153349) (0.104719755,0.973369646) (0.125663706,0.971190671) (0.146607657,0.96861709) (0.167551608,0.965649687) (0.188495559,0.962289369) (0.20943951,0.958537162) (0.230383461,0.954394214) (0.251327412,0.949861794) (0.272271363,0.94494129) (0.293215314,0.939634213) (0.314159265,0.933942194) (0.335103216,0.927866983) (0.356047167,0.921410455) (0.376991118,0.914574601) (0.397935069,0.907361536) (0.41887902,0.899773496) (0.439822972,0.891812834) (0.460766923,0.883482029) (0.481710874,0.874783677) (0.502654825,0.865720497) (0.523598776,0.856295327) (0.544542727,0.846511126) (0.565486678,0.836370977) (0.586430629,0.825878079) (0.60737458,0.815035754) (0.628318531,0.803847447) (0.649262482,0.79231672) (0.670206433,0.780447258) (0.691150384,0.768242867) (0.712094335,0.755707472) (0.733038286,0.74284512) (0.753982237,0.729659979) (0.774926188,0.716156338) (0.795870139,0.702338606) (0.81681409,0.688211314) (0.837758041,0.673779111) (0.858701992,0.65904677) (0.879645943,0.644019183) (0.900589894,0.628701364) (0.921533845,0.613098447) (0.942477796,0.597215687) (0.963421747,0.581058459) (0.984365698,0.564632261) (1.00530965,0.547942709) (1.0262536,0.530995542) (1.04719755,0.513796618) (1.0681415,0.496351918) (1.08908545,0.478667542) (1.1100294,0.460749712) (1.13097336,0.442604769) (1.15191731,0.424239177) (1.17286126,0.40565952) (1.19380521,0.386872501) (1.21474916,0.367884947) (1.23569311,0.348703804) (1.25663706,0.329336139) (1.27758101,0.309789139) (1.29852496,0.290070113) (1.31946891,0.270186491) (1.34041287,0.250145823) (1.36135682,0.22995578) (1.38230077,0.209624153) (1.40324472,0.189158856) (1.42418867,0.168567922) (1.44513262,0.147859504) (1.46607657,0.127041879) (1.48702052,0.106123442) (1.50796447,0.0851127097) (1.52890842,0.0640183195) (1.54985238,0.0428490296) (1.57079633,0.0216137193) (1.59174028,0.000321388316) (1.61268423,-0.0210188426) (1.63362818,-0.042397732) (1.65457213,-0.0638059175) (1.67551608,-0.0852339161) (1.69646003,-0.106672124) (1.71740398,-0.128110816) (1.73834793,-0.149540147) (1.75929189,-0.170950151) (1.78023584,-0.19233074) (1.80117979,-0.213671707) (1.82212374,-0.234962723) (1.84306769,-0.256193338) (1.86401164,-0.277352983) (1.88495559,-0.298430965) (1.90589954,-0.319416473) (1.92684349,-0.340298575) (1.94778745,-0.361066216) (1.9687314,-0.381708222) (1.98967535,-0.402213299) (2.0106193,-0.422570029) (2.03156325,-0.442766878) (2.0525072,-0.462792186) (2.07345115,-0.482634176) (2.0943951,-0.502280948) (2.11533905,-0.521720482) (2.136283,-0.540940639) (2.15722696,-0.559929155) (2.17817091,-0.578673649) (2.19911486,-0.597161617) (2.22005881,-0.615380436) (2.24100276,-0.63331736) (2.26194671,-0.650959525) (2.28289066,-0.668293943) (2.30383461,-0.685307507) (2.32477856,-0.701986991) (2.34572251,-0.718319043) (2.36666647,-0.734290196) (2.38761042,-0.749886858) (2.40855437,-0.765095319) (2.42949832,-0.779901746) (2.45044227,-0.794292187) (2.47138622,-0.808252567) (2.49233017,-0.821768693) (2.51327412,-0.834826249) (2.53421807,-0.847410799) (2.55516202,-0.859507786) (2.57610598,-0.871102533) (2.59704993,-0.882180242) (2.61799388,-0.892725992) (2.63893783,-0.902724744) (2.65988178,-0.912161336) (2.68082573,-0.921020488) (2.70176968,-0.929286797) (2.72271363,-0.93694474) (2.74365758,-0.943978672) (2.76460154,-0.950372828) (2.78554549,-0.956111323) (2.80648944,-0.961178151) (2.82743339,-0.965557183) (2.84837734,-0.969232173) (2.86932129,-0.97218675) (2.89026524,-0.974404426) (2.91120919,-0.97586859) (2.93215314,-0.97656251) (2.95309709,-0.976469334) (2.97404105,-0.97557209) (2.994985,-0.973853684) (3.01592895,-0.971296901) (3.0368729,-0.967884405) (3.05781685,-0.963598741) (3.0787608,-0.958422332) (3.09970475,-0.95233748) (3.1206487,-0.945326366) (3.14159265,-0.937371052)};
\addlegendentry{正交投影}
\addplot+[mark=none] coordinates{(-3.14159265,0.123909926) (-3.1206487,0.0823331616) (-3.09970475,0.0424536449) (-3.0787608,0.00424280261) (-3.05781685,-0.0323277463) (-3.0368729,-0.0672861903) (-3.01592895,-0.100660525) (-2.994985,-0.132478555) (-2.97404105,-0.162767892) (-2.95309709,-0.191555953) (-2.93215314,-0.218869965) (-2.91120919,-0.244736963) (-2.89026524,-0.269183789) (-2.86932129,-0.29223709) (-2.84837734,-0.313923324) (-2.82743339,-0.334268756) (-2.80648944,-0.353299458) (-2.78554549,-0.371041308) (-2.76460154,-0.387519994) (-2.74365758,-0.402761012) (-2.72271363,-0.416789662) (-2.70176968,-0.429631055) (-2.68082573,-0.441310109) (-2.65988178,-0.451851549) (-2.63893783,-0.461279906) (-2.61799388,-0.469619523) (-2.59704993,-0.476894545) (-2.57610598,-0.48312893) (-2.55516202,-0.488346439) (-2.53421807,-0.492570643) (-2.51327412,-0.495824921) (-2.49233017,-0.498132459) (-2.47138622,-0.499516249) (-2.45044227,-0.499999092) (-2.42949832,-0.499603598) (-2.40855437,-0.498352182) (-2.38761042,-0.496267068) (-2.36666647,-0.493370287) (-2.34572251,-0.489683679) (-2.32477856,-0.485228889) (-2.30383461,-0.480027371) (-2.28289066,-0.474100388) (-2.26194671,-0.467469008) (-2.24100276,-0.460154109) (-2.22005881,-0.452176375) (-2.19911486,-0.443556297) (-2.17817091,-0.434314175) (-2.15722696,-0.424470117) (-2.136283,-0.414044037) (-2.11533905,-0.403055658) (-2.0943951,-0.391524509) (-2.07345115,-0.379469928) (-2.0525072,-0.36691106) (-2.03156325,-0.353866857) (-2.0106193,-0.340356081) (-1.98967535,-0.326397299) (-1.9687314,-0.312008886) (-1.94778745,-0.297209025) (-1.92684349,-0.282015707) (-1.90589954,-0.26644673) (-1.88495559,-0.250519701) (-1.86401164,-0.234252031) (-1.84306769,-0.217660943) (-1.82212374,-0.200763465) (-1.80117979,-0.183576433) (-1.78023584,-0.16611649) (-1.75929189,-0.148400088) (-1.73834793,-0.130443487) (-1.71740398,-0.112262751) (-1.69646003,-0.0938737567) (-1.67551608,-0.0752921841) (-1.65457213,-0.0565335228) (-1.63362818,-0.0376130698) (-1.61268423,-0.0185459294) (-1.59174028,0.000652986295) (-1.57079633,0.0199689578) (-1.54985238,0.0393874578) (-1.52890842,0.0588941516) (-1.50796447,0.0784748969) (-1.48702052,0.0981157437) (-1.46607657,0.117802934) (-1.44513262,0.137522904) (-1.42418867,0.15726228) (-1.40324472,0.177007881) (-1.38230077,0.196746721) (-1.36135682,0.216466003) (-1.34041287,0.236153124) (-1.31946891,0.255795675) (-1.29852496,0.275381436) (-1.27758101,0.294898381) (-1.25663706,0.314334678) (-1.23569311,0.333678686) (-1.21474916,0.352918955) (-1.19380521,0.37204423) (-1.17286126,0.391043447) (-1.15191731,0.409905734) (-1.13097336,0.428620413) (-1.1100294,0.447176997) (-1.08908545,0.465565192) (-1.0681415,0.483774897) (-1.04719755,0.501796202) (-1.0262536,0.519619389) (-1.00530965,0.537234936) (-0.984365698,0.554633509) (-0.963421747,0.57180597) (-0.942477796,0.58874337) (-0.921533845,0.605436956) (-0.900589894,0.621878164) (-0.879645943,0.638058625) (-0.858701992,0.653970161) (-0.837758041,0.669604788) (-0.81681409,0.684954712) (-0.795870139,0.700012332) (-0.774926188,0.714770242) (-0.753982237,0.729221226) (-0.733038286,0.74335826) (-0.712094335,0.757174514) (-0.691150384,0.77066335) (-0.670206433,0.783818322) (-0.649262482,0.796633176) (-0.628318531,0.809101851) (-0.60737458,0.82121848) (-0.586430629,0.832977385) (-0.565486678,0.844373082) (-0.544542727,0.855400282) (-0.523598776,0.866053883) (-0.502654825,0.876328981) (-0.481710874,0.886220861) (-0.460766923,0.895725001) (-0.439822972,0.904837072) (-0.41887902,0.913552937) (-0.397935069,0.921868651) (-0.376991118,0.929780463) (-0.356047167,0.937284813) (-0.335103216,0.944378333) (-0.314159265,0.951057849) (-0.293215314,0.957320379) (-0.272271363,0.963163132) (-0.251327412,0.968583511) (-0.230383461,0.97357911) (-0.20943951,0.978147718) (-0.188495559,0.982287313) (-0.167551608,0.985996068) (-0.146607657,0.989272347) (-0.125663706,0.992114707) (-0.104719755,0.994521897) (-0.0837758041,0.99649286) (-0.0628318531,0.998026729) (-0.041887902,0.99912283) (-0.020943951,0.999780683) (-4.4408921e-16,1) (0.020943951,0.999780683) (0.041887902,0.99912283) (0.0628318531,0.998026729) (0.0837758041,0.99649286) (0.104719755,0.994521897) (0.125663706,0.992114707) (0.146607657,0.989272347) (0.167551608,0.985996068) (0.188495559,0.982287313) (0.20943951,0.978147718) (0.230383461,0.97357911) (0.251327412,0.968583511) (0.272271363,0.963163132) (0.293215314,0.957320379) (0.314159265,0.951057849) (0.335103216,0.944378333) (0.356047167,0.937284813) (0.376991118,0.929780463) (0.397935069,0.921868651) (0.41887902,0.913552937) (0.439822972,0.904837072) (0.460766923,0.895725001) (0.481710874,0.886220861) (0.502654825,0.876328981) (0.523598776,0.866053883) (0.544542727,0.855400282) (0.565486678,0.844373082) (0.586430629,0.832977385) (0.60737458,0.82121848) (0.628318531,0.809101851) (0.649262482,0.796633176) (0.670206433,0.783818322) (0.691150384,0.77066335) (0.712094335,0.757174514) (0.733038286,0.74335826) (0.753982237,0.729221226) (0.774926188,0.714770242) (0.795870139,0.700012332) (0.81681409,0.684954712) (0.837758041,0.669604788) (0.858701992,0.653970161) (0.879645943,0.638058625) (0.900589894,0.621878164) (0.921533845,0.605436956) (0.942477796,0.58874337) (0.963421747,0.57180597) (0.984365698,0.554633509) (1.00530965,0.537234936) (1.0262536,0.519619389) (1.04719755,0.501796202) (1.0681415,0.483774897) (1.08908545,0.465565192) (1.1100294,0.447176997) (1.13097336,0.428620413) (1.15191731,0.409905734) (1.17286126,0.391043447) (1.19380521,0.37204423) (1.21474916,0.352918955) (1.23569311,0.333678686) (1.25663706,0.314334678) (1.27758101,0.294898381) (1.29852496,0.275381436) (1.31946891,0.255795675) (1.34041287,0.236153124) (1.36135682,0.216466003) (1.38230077,0.196746721) (1.40324472,0.177007881) (1.42418867,0.15726228) (1.44513262,0.137522904) (1.46607657,0.117802934) (1.48702052,0.0981157437) (1.50796447,0.0784748969) (1.52890842,0.0588941516) (1.54985238,0.0393874578) (1.57079633,0.0199689578) (1.59174028,0.000652986295) (1.61268423,-0.0185459294) (1.63362818,-0.0376130698) (1.65457213,-0.0565335228) (1.67551608,-0.0752921841) (1.69646003,-0.0938737567) (1.71740398,-0.112262751) (1.73834793,-0.130443487) (1.75929189,-0.148400088) (1.78023584,-0.16611649) (1.80117979,-0.183576433) (1.82212374,-0.200763465) (1.84306769,-0.217660943) (1.86401164,-0.234252031) (1.88495559,-0.250519701) (1.90589954,-0.26644673) (1.92684349,-0.282015707) (1.94778745,-0.297209025) (1.9687314,-0.312008886) (1.98967535,-0.326397299) (2.0106193,-0.340356081) (2.03156325,-0.353866857) (2.0525072,-0.36691106) (2.07345115,-0.379469928) (2.0943951,-0.391524509) (2.11533905,-0.403055658) (2.136283,-0.414044037) (2.15722696,-0.424470117) (2.17817091,-0.434314175) (2.19911486,-0.443556297) (2.22005881,-0.452176375) (2.24100276,-0.460154109) (2.26194671,-0.467469008) (2.28289066,-0.474100388) (2.30383461,-0.480027371) (2.32477856,-0.485228889) (2.34572251,-0.489683679) (2.36666647,-0.493370287) (2.38761042,-0.496267068) (2.40855437,-0.498352182) (2.42949832,-0.499603598) (2.45044227,-0.499999092) (2.47138622,-0.499516249) (2.49233017,-0.498132459) (2.51327412,-0.495824921) (2.53421807,-0.492570643) (2.55516202,-0.488346439) (2.57610598,-0.48312893) (2.59704993,-0.476894545) (2.61799388,-0.469619523) (2.63893783,-0.461279906) (2.65988178,-0.451851549) (2.68082573,-0.441310109) (2.70176968,-0.429631055) (2.72271363,-0.416789662) (2.74365758,-0.402761012) (2.76460154,-0.387519994) (2.78554549,-0.371041308) (2.80648944,-0.353299458) (2.82743339,-0.334268756) (2.84837734,-0.313923324) (2.86932129,-0.29223709) (2.89026524,-0.269183789) (2.91120919,-0.244736963) (2.93215314,-0.218869965) (2.95309709,-0.191555953) (2.97404105,-0.162767892) (2.994985,-0.132478555) (3.01592895,-0.100660525) (3.0368729,-0.0672861903) (3.05781685,-0.0323277463) (3.0787608,0.00424280261) (3.09970475,0.0424536449) (3.1206487,0.0823331616) (3.14159265,0.123909926)};
\addlegendentry{Taylor}
\end{axis}\end{tikzpicture}\end{center}

\end{solution}
